\documentclass[final]{beamer}

\pdfvariable suppressoptionalinfo 512\relax

% ====================
% Packages
% ====================

\usepackage[T1]{fontenc}
\usepackage{../hiromacros}

\usepackage{lmodern}
\usepackage[size=A0,orientation=portrait,scale=1.0]{beamerposter}
\usetheme{gemini}
\usecolortheme{mit}
\usepackage{graphicx}
\usepackage{booktabs}
\usepackage{tikz}
\usepackage{pgfplots}
\pgfplotsset{compat=1.14}
\usepackage{physics}
\usepackage{cleveref}
\usepackage[font=scriptsize]{caption}
\usepackage[list=true, font=scriptsize, labelformat=brace]{subcaption}
\graphicspath{{figs/}}

% Bibliographic Stuff
\usepackage[backend=biber, language=english, style=phys]{biblatex}
\addbibresource{references.bib}

\DeclareDocumentCommand\vectorbold{ s m
}{\IfBooleanTF{#1}{\symbfit{#2}}{\symbf{#2}}} % Vector bold [star for
% Greek and italic Roman]

\makeatletter
\newcommand\thefontsize[1]{{#1 The current font size is: \f@size pt\par}}
\makeatother

% ====================
% Lengths
% ====================

% If you have N columns, choose \sepwidth and \colwidth such that
% (N+1)*\sepwidth + N*\colwidth = \paperwidth
\newlength{\sepwidth}
\newlength{\colwidth}
\setlength{\sepwidth}{0.066\paperwidth}
\setlength{\colwidth}{0.4\paperwidth}

\newcommand{\separatorcolumn}{\begin{column}{\sepwidth}\end{column}}

% ====================
% Title
% ====================

\title{Calculating Energy Flows in Strongly Coupled Open Quantum
  Systems with HOPS}

\author{\underline{Valentin Boettcher}\inst{1}, Richard Hartmann\inst{1},
  Konstantin Beyer\inst{1}, Walter Strunz\inst{1}}

\institute[shortinst]{\inst{1} Institute for Theoretical Physics, Dresden, Germany}

% ====================
% Footer (optional)
% ====================

\footercontent{
  \href{https://tu-dresden.de/mn/physik/itp/tqo/die-professur}{https://tu-dresden.de/mn/physik/itp/tqo/die-professur} \hfill
  Group for Theoretical Quantum Optics, TU-Dresden \hfill
  \href{mailto:valentin.boettcher@tu-dresden.de}{valentin.boettcher@tu-dresden.de}}
% (can be left out to remove footer)

% ====================
% Logo (optional)
% ====================

% use this to include logos on the left and/or right side of the header:
\logoleft{\includegraphics[height=3cm]{Logo_TU_Dresden.pdf}}

% ====================
% Body
% ====================

\begin{document}

\begin{frame}[t]
\begin{columns}[t]
\separatorcolumn

\begin{column}{\colwidth}
  \begin{block}{Premise}
    \begin{itemize}
    \item application of thermodynamic notions to strongly coupled and
      non-Markovian open quantum systems is non-trivial
      (\cite{Rivas2019Oct,Kato2016Dec,Strasberg2021Aug,Talkner2020Oct} and many
      more)
    \item dynamics of bath and interaction Hamiltonians plays an
      important role \(\rightarrow\) must not be neglected in the
      strong coupling regime
    \item the ``Hierarchy of Pure States''
      (HOPS~\cite{Hartmann2017Dec,Diosi1998Mar}) gives the us ability to
      simulate open quantum systems \textbf{exactly}
    \item because HOPS simulates global dynamics \(\rightarrow\) \textbf{gives
      access to certain bath dynamics with no additional effort}
    \end{itemize}
  \end{block}
  \begin{block}{NMQSD/HOPS}
    \begin{itemize}
    \item consider the model of a general quantum system (\(H_\sys(t)\))
    coupled to \(N\) baths
    \begin{equation}
      \label{eq:generalmodel}
      H(t) = H_\sys(t) + ∑_{n=1}^N \qty[L_n^†(t)B_n + \hc] + ∑_{n=1}^NH_B\nth ,
    \end{equation}
    with \(B_n=∑_{λ} g_λ\nth a_λ\nth\) and
    \(H_B\nth=∑_λω_λ\nth \qty(b_λ\nth)^\dag b_λ\nth\)

    \item leads to \emph{stochastic} Non-Markovian
    Quantum State Diffusion (NMQSD)
    \begin{equation}
      \label{eq:nmqsd}
      ∂_tψ_t(\vb{η}^\ast_t) = -\iu H(t) ψ_t(\vb{η}^\ast_t) +
      \vb{L}\cdot \vb{η}^\ast_tψ_t(\vb{η}^\ast_t) - ∑_{n=1}^N L(t)_n^†∫_0^t\dd{s}α_n(t-s)\fdv{ψ_t(\vb{η}^\ast_t)}{η^\ast_n(s)}
    \end{equation}
    \item the reduced state of the system is recovered through
      \(ρ=\mathcal{M}(ψ_t(\vb{η}^\ast_t)ψ_t^\dag(\vb{η}^\ast_t))\)

    \item with \(α_n(τ)=∑_{\mu}^{M_n}G_μ\nth\eu^{-W_μ\nth τ}\)
      we find for the hierarchy state \(\ket{Ψ}\in \hilb\otimes\mathcal{F}\)
    \begin{equation}
      \label{eq:fockhops}
      \begin{aligned}
        ∂_t\ket{Ψ} &= \qty[
                     \begin{aligned}
                       -\iu H_\sys + \vb{L}\cdot\vb{η}^\ast &-
                                                              ∑_{n=1}^N∑_{μ=1}^{M_n}b_{n,μ}^\dag b_{n,μ} W\nth_μ \\
                                                            &\qquad+
                                                              \iu ∑_{n=1}^N∑_{μ=1}^{M_n} \sqrt{G_{n,μ}} \qty(b^†_{n,μ}L_n +
                                                              b_{n,μ}L^†_n)
                     \end{aligned}
                     ] \ket{Ψ}.
      \end{aligned}
    \end{equation}
    \item truncating hierarchy depth \(\kmat\) in \cref{eq:fockhops}
      yields numeric method

    \item finite temperature \(\rightarrow\) substitute
    \(B(t)\rightarrow B(t)+ξ(t)\)
    \item \(\exists\) nonlinear method which improves convergence drastically
    \end{itemize}

    See~\cite{Hartmann2017Dec} for details.
  \end{block}

  \begin{alertblock}{Bath Observables}
    \begin{itemize}
    \item \cref{eq:nmqsd,eq:fockhops} \(\implies\) correspondence
      \(B^n(t) \leftrightarrow ψ^\kmat\) {\tiny(\(\abs{\kmat}=n\))}

    \item can calulate observables of type
    \(O_\sys\otimes (B^a)^\dag B^b\) + time derivatives
    \end{itemize}

    We consider the zero temperature case with one bath in the linear
    method.

    \begin{description}
    \item[Bath Energy Flow]
      \begin{equation}
        \label{eq:heatflowdef}
        \begin{aligned}
          J &= - \dv{\ev{H_\bath}}{t}  = \ev{L(t)^†∂_t B(t) + L(t)∂_t
              B^†(t)}_\inter \\
            &=-\i \mathcal{M}_{η^\ast}\bra{\psi(η,
              t)}L(t)^†\dot{D}_t\ket{\psi(η^\ast,t)} + \cc\\
            &= - ∑_\mu\sqrt{G_\mu}W_\mu
              \mathcal{M}_{η^\ast}\bra{\psi^{(0)}(η,t)}L(t)^†\ket{\psi^{\vb{e}_\mu}(η^\ast,t)} + \cc
        \end{aligned}
      \end{equation}
      The expectation value of bath energy flow is connected to the first
      level hierarchy states.
    \item[Interaction Energy]
      A similar expression exists for the expectation value of the
        interaction energy.
      \begin{equation}
        \label{eq:intexp}
        \begin{aligned}
          \ev{H_\inter} &=-\i \mathcal{M}_{η^\ast}\bra{\psi(η,
          t)}L(t)^†D_t\ket{\psi(η^\ast,t)} + \cc \\
                        &=  ∑_\mu\sqrt{G_\mu}
                          \mathcal{M}_{η^\ast}\bra{\psi^{(0)}(η,t)}L(t)^†\ket{\psi^{\vb{e}_\mu}(η^\ast,t)} + \cc.
        \end{aligned}
      \end{equation}
    \end{description}

    This result allows us (in principle) to calculate the energy flow in
    \textbf{arbitrarily driven systems} for a \textbf{wide temperature
      range} and with (sub-)Ohmic BCF.
  \end{alertblock}

  \begin{block}{Comparison to an Analytic Solution}
    \begin{columns}
      \begin{column}{.5\colwidth}
        \begin{itemize}
        \item ``Quantum Brownian Motion'' like model
          \begin{equation}
            \label{eq:hamiltonian}
            \begin{aligned}
              H = ∑_{i\in\qty{1,2}} &\qty[H^{(i)}_O + q_iB^{(i)} +
                                      H_B^{(i)}]\\
                                    &\quad+ \frac{γ}{4}(q_1-q_2)^2
            \end{aligned}
          \end{equation}
          where \(H_O^{(i)}= \frac{Ω_i}{4}\qty(p_i^2+q_i^2)\)
        \item exact solution via exponential expansion of the BCF in the
          Heisenberg picture \(\rightarrow\) easy access to bath energy
          flow
        \item here \tval{analytic/omega}, \tval{analytic/gamma},
          \(α(τ)=η (1+\iu ω_cτ)^{-(2)}\) with
          \tval{analytic/cutoff_freq}, \tval{analytic/bcf_zero}
        \end{itemize}
      \end{column}
      \begin{column}{.5\colwidth}
        \begin{figure}[H]
          \centering
          \plot{analytic/flow}
          \caption{\label{fig:brownian}The bath energy flows \cref{eq:heatflowdef} for the
            hot (lower line) and the cold (upper line) bath show that the
            analytical and numerical results are compatible.}
        \end{figure}
      \end{column}
    \end{columns}
  \end{block}
  \begin{block}{Possible Applications}
    \begin{itemize}
    \item \emph{Simulation of thermal quantum machines}
    \item convergence criteria for HOPS: energy conservation, calculating the
      same observable in multiple ways
    \item quantification of entanglement of system and bath
    \item testing results obtained from approximations
    \end{itemize}
  \end{block}
  \begin{block}{Resources}
    {\AtNextBibliography{\tiny} \printbibliography}
  \end{block}
\end{column}

\separatorcolumn

\begin{column}{\colwidth}
  \begin{block}{Spin-Boson like Model and BCF Dependence}
    \begin{itemize}
    \item spin-boson like model coupled to a zero temperature bath
      \begin{equation}
        \label{eq:spinbos}
        H_\sys= \frac12 σ_z,\, L=\frac12 σ_x,\, α(τ)=η (1+\iu ω_cτ)^{-(2)}
      \end{equation}
    \item memory time \(\sim 1/ω_c\) has qualitative influence on the
      bath energy flow
    \end{itemize}
    \begin{figure}[H]
      \centering
      \begin{subfigure}[t]{.49\columnwidth}
        \plot{one_bath/omega_interaction}
        \caption{\label{fig:omega_ints}\tval{one_bath/omega_bcf_str}}
      \end{subfigure}
      \begin{subfigure}[t]{.49\columnwidth}
        \plot{one_bath/delta_interaction}
        \caption{\tval{one_bath/delta_bcf_wc}}
      \end{subfigure}
      \caption{The interaction energy expectation value for different
        cutoff frequencies and coupling strengths, where the dashed
        lines are obtained using energy conservation while the solid
        lines are the result of direct calculation. The percentages in
        the legend tell how many points are compatible within one
        standard deviation. The Statistical error estimate is smaller
        than the line width. \(N=5\cdot 10^5\) trajectories have been
        used.}
    \end{figure}
  \end{block}
  \begin{block}{Initial Slip}
    \begin{itemize}
    \item for \emph{very} short times \(\rightarrow\) \(H_\sys\approx
      0\), origin of the \emph{``Initial Slip''} spike in
      \cref{fig:brownian}:
      \begin{equation}
        \label{eq:purede}
        \ev{\dot{H}_\bath } = -2 ∫_0^t\dd{s}\ev{L(t)L(s)} \Im[\dot{α}(t-s)].
      \end{equation}
    \item determines ultra short-time shape of \emph{all} trajectories
    \end{itemize}
    \begin{figure}[H]
      \centering
      \begin{subfigure}[t]{.49\linewidth}
        \plot{one_bath/initial_slip}
        \caption{\label{fig:initslipconst}The bath energy flows for the same settings as in
          \cref{fig:omega_ints}, where the dashed lines correspond to \cref{eq:purede}.}
      \end{subfigure}
      \begin{subfigure}[t]{.49\linewidth}
        \plot{modcoup/initial_slip_modcoup}
        \caption{Same as \cref{fig:initslipconst}, but for modulated
          coupling (``Smoothstep'' functions with smoothness \(s\), see inset).}
      \end{subfigure}
    \end{figure}
  \end{block}
  \begin{block}{Modulating the Coupling}
    \begin{itemize}
    \item same model as above \cref{eq:spinbos}, but with \(L(τ) =
      \sin^2(\frac{Δ}{2} τ)σ_x\)
    \item Question: How much energy can be extracted from a system
      connected to a single bath? (Ergotropy)
      \begin{itemize}
      \item Answer: less than
        \(ΔE_{\mathrm{max}}=\frac{1}{β}\qrelent{ρ}{ρ_β}\)
      \end{itemize}
    \end{itemize}
    \begin{figure}[H]
      \centering
      \begin{subfigure}[t]{.49\linewidth}
        \plot{modcoup/omegas_total}
        \caption{\label{fig:omega_total}The total energy for
          \tval{modcoup/omega_delta} and \tval{modcoup/omega_alpha}
          but varying cutoff. Energy is normalized to the
          ergotropy. The dashed vertical lines illustrate the bath
          memory time (\(|α(τ)| = α(0)/300\)).}
      \end{subfigure}
      % \begin{subfigure}[t]{.49\linewidth}
      %   \plot{modcoup/flow_interaction_overview}
      %   \caption{Same situation as in \cref{fig:omega_total}. All
      %     quantities are normalized to the \(ω_c=1\)
      %     case.}
      % \end{subfigure}
      \begin{subfigure}[t]{.49\linewidth}
        \plot{modcoup/delta_dependence}
        \caption{Maximum one-shot power for \(10\) periods and
          different modulation frequencies \(Δ\) and coupling
          strengths.}
      \end{subfigure}
    \end{figure}
  \end{block}
  \begin{block}{Continuously Coupled Engine (Preliminary)}
    \begin{itemize}
    \item qubit coupled to two baths of different
      temperatures (\(T_c, T_h\))
      \begin{equation}
        \label{eq:antizenomodel}
        H_\sys= \frac12  \qty[ω_0 + γ Δ\sin(Δ t)]σ_z,\, L_{c,h}=\frac12 σ_x
      \end{equation}
      {\tiny \tval{anti_zeno/delta}, \tval{anti_zeno/gamma}, \tval{anti_zeno/omega_alpha},
        \tval{anti_zeno/omega_zero}, \tval{anti_zeno/tc}, \tval{anti_zeno/th}}
    \item system Hamiltonian modulated + baths periodically
      decoupled (cooldown) and ``reset''
    \end{itemize}
    \begin{figure}[H]
      \centering
      \begin{subfigure}[t]{.49\linewidth}
        \plot{anti_zeno/sd_setup}
        \caption{The spectral densities of the baths, where the vertical
          lines show where \(ω=ω_0 \pm Δ\). The overlap of the filter
          for \(n\) modulation periods is crucial for the ``anti-zeno''
          effect~\cite{Mukherjee2020Jan}.}
      \end{subfigure}
      % \begin{subfigure}[t]{.49\linewidth}
      %   \plot{modcoup/flow_interaction_overview}
      %   \caption{Same situation as in \cref{fig:omega_total}. All
      %     quantities are normalized to the \(ω_c=1\)
      %     case.}
      % \end{subfigure}
      \begin{subfigure}[t]{.49\linewidth}
        \plot{anti_zeno/modulation_setup}
        \caption{Setup of the system and coupling modulation for the
          case with ``cooldown''. The units are arbitrary. The
          initialization period and two cycles are shown.}
      \end{subfigure}
      \begin{subfigure}[t]{.49\linewidth}
        \plot{anti_zeno/anti_zeno_with_cool}
        \caption{\label{fig:cont_coup} The total energy change after the initialization
          period, where vertical lines show the points at which the
          times for the calculation of the power are taken. The mean
          power obtained is \tval{anti_zeno/power_with_cool}.}.
      \end{subfigure}
      \begin{subfigure}[t]{.49\linewidth}
        \plot{anti_zeno/anti_zeno_without_cool}
        \caption{Same as in \cref{fig:cont_coup} but without
          cooldown. The mean power obtained is
          \tval{anti_zeno/power_without_cool}. The difference is not
          significant yet.}.
      \end{subfigure}
    \end{figure}
  \end{block}
\end{column}

\separatorcolumn
\end{columns}
\end{frame}

\end{document}
